% 原书第 88—92 页；第 6.6.2—6.6.3 节。
\subsection{预测模拟}
\label{subsec:ch06-prediction-simulation}

我们评价了 HAL 估计量相对于其他非参数算法的有限样本表现。比较算法包括回归树（regression tree；Breiman et al.，1984）、随机森林（random forest；Breiman，2001）、梯度提升机（gradient boosted machine，GBM；Friedman，2001）、核回归（kernel regression；Nadaraya，1964；Watson，1964）、支持向量机（support vector machine，SVM；Hearst et al.，1998），以及多项式多元自适应回归样条（Polynomial MARS；Friedman，1991）。我们考虑了三类数据生成机制，分别称为光滑型、跳跃型和正弦型。对于每一类数据生成机制，改变 $W$ 的维数，取 $d\in\{1,3,5\}$，样本量取 $n\in\{500,1000,2000\}$。以 $R^2$ 评价表现，并在一个样本量为 $N=1\mathrm e4$ 的独立测试集上计算。对于给定估计量 $\bar Q_n$，定义
\[
 R^2=1-\frac{\sum_{i=1}^N\{Y_i-\bar Q_n(W_i)\}^2}
                  {\sum_{i=1}^N\{Y_i-\bar Y_N\}^2}.
\]
每种情形都设计为使最优的 $R^2$ 值等于 $R^2_{opt}=0.80$，其中，
\[
 R^2_{opt}=1-\frac{E_{P_0}\{Y-\bar Q_0(W)\}^2}{\Var_0(Y)}
\]
是使用真实回归函数 $\bar Q_0$ 时得到的 $R^2$ 值。这个值可以看作任何估计量表现的上界。

$W$ 的分布如下：$W_1\sim\operatorname{Uniform}(-4,4)$，$W_2\sim\operatorname{Uniform}(-4,4)$，$W_3\sim\operatorname{Bernoulli}(0.5)$，$W_4\sim\operatorname{Normal}(0,1)$，$W_5\sim\operatorname{Gamma}(2,1)$。对于维数 $d$，将目标参数记为 $\bar Q_0^d(W)$，令 $W=(W_j:j=1,\ldots,d\})$。定义 $Y=\bar Q_0^d(W)+\epsilon$，其中，$\epsilon\sim\operatorname{Normal}(0,1)$。
% 上式 W 的右端多一个右花括号，按原书印刷保留。

对于 $d=1,3,5$，“光滑型”回归函数分别定义为
\[
 \begin{aligned}
 \bar Q_0^1(w)&=0.05w_1+0.42w_1^2;\\
 \bar Q_0^3(w)&=0.07w_1-0.28w_1^2+0.5w_2+0.25w_2w_3;\\
 \bar Q_0^5(w)&=0.1w_1-0.3w_1^2+0.25w_2+0.5w_3w_2
                         -0.5w_4+0.04w_5^2-0.1w_5.
 \end{aligned}
\]

“跳跃型”回归函数定义为
\[
 \begin{aligned}
 \bar Q_0^1(w)={}&-2.7I(w_1<-3)+2.5I(w_1>-2)-2I(w_1>0)\\
                &{}+4I(w_1>2)-3I(w_1>3);\\[.4em]
 \bar Q_0^3(w)={}&-2I(w_1<-3)w_3+2.5I(w_1>-2)-2I(w_1>0)\\
                &{}+2.5I(w_1>2)w_3-2.5I(w_1>3)+I(w_2>-1)\\
                &{}-4I(w_2>1)w_3+2I(w_2>3);\\[.4em]
 \bar Q_0^5(w)={}&-I(w_1<-3)w_3+0.5I(w_1>-2)-I(w_1>0)\\
                &{}+2I(w_1>2)w_3-3I(w_1>3)+1.5I(w_2>-1)\\
                &{}-5I(w_2>1)w_3+2I(w_2>3)+2I(w_4<0)\\
                &{}-I(w_5>5)-I(w_4<0)I(w_1<0)+2w_3.
 \end{aligned}
\]

“正弦型”回归函数定义为
\[
 \begin{aligned}
 \bar Q_0^1(w)={}&2\sin(0.5\pi|w_1|)+2\cos(0.5\pi|w_1|);\\
 \bar Q_0^3(w)={}&4w_3I(w_2<0)\sin(0.5\pi|w_1|)
                    +4.1I(w_2\geq0)\cos(0.5\pi|w_1|);\\
 \bar Q_0^5(w)={}&3.8w_3I(w_2<0)\sin(0.5\pi|w_1|)
                    +4I(w_2>0)\cos(\pi|w_1|/2)\\
                &{}+0.1w_5\sin(\pi w_4)+w_3\cos(|w_4-w_5|).
 \end{aligned}
\]

图 6.3 展示了模拟研究的结果，其中，行表示不同的数据生成机制，列表示 $W$ 的不同维数。图的边缘面板给出了汇总结果：底部边缘面板在给定维数下，汇总不同数据生成机制的结果；右侧边缘面板则在给定数据生成机制下，汇总不同维数的结果。每个面板中的算法，都按三个样本量下的平均 $R^2$ 值排序：$R^2$ 最高的在顶部，最低的在底部。

先看对应“光滑型”数据生成机制的第一行。当 $d=1$ 时，除随机森林外，所有算法都表现良好，其中核回归最好。但随着维数增加，核回归的相对表现下降，而 HAL 的相对表现上升。综合所有维数，SVM 的总体表现最好，不过 GBM 和 HAL 的表现也与之相近。在对应“跳跃型”情形的第二行中，HAL 表现极佳；对于所有维数，当 $n=2000$ 时，它的 $R^2$ 几乎达到了最优值。在对应“正弦型”情形的第三行中，略令人意外的是，核回归在所有维数下都表现最好。这似乎部分是由于，在 $n=500$ 时，它优于其他估计量。样本量较大时，核回归、随机森林和 HAL 达到的 $R^2$ 都同样很高。右下角最远处的面板汇总了所有模拟的结果，并按平均 $R^2$ 对算法排序；可以看到，HAL 的平均 $R^2$ 最高，其次是核回归和随机森林。总体而言，HAL 在所有情形中，相对于其他算法都表现良好；在局部光滑性假设不成立的跳跃型情形中，它的表现尤其突出。

\begin{landscape}
\begin{figure}[p]
 \centering
 % 图 6.3：原始矢量数据层原样保留；仅遮盖文字区，按原顺序覆排中文标签。
\resizebox{.94\linewidth}{!}{%
\begin{tikzpicture}[x=1bp,y=-1bp]
 \node[anchor=north west,inner sep=0pt] at (0,0) {\includegraphics[width=475bp,height=282bp]{figures/ch06_fig3_source.pdf}};
 \fill[white] (0,0) rectangle (35,282);
 \fill[white] (160,0) rectangle (250,15);
 \fill[white] (437,17) rectangle (454,26.7);
 \node[font=\fontsize{10.5}{11}\selectfont] at (204,7.2) {维数};
 \node[rotate=90,font=\fontsize{10.5}{11}\selectfont] at (10.5,112) {情景};
 \node[rotate=90,font=\fontsize{7.5}{8}\selectfont] at (28,47) {光滑};
 \node[rotate=90,font=\fontsize{7.5}{8}\selectfont] at (28,113) {跳跃};
 \node[rotate=90,font=\fontsize{7.5}{8}\selectfont] at (28,178) {正弦};
 \node[rotate=90,font=\fontsize{7.5}{8}\selectfont] at (28,244) {全部};
 \node[font=\fontsize{7.5}{8}\selectfont] at (445,23.5) {全部};
 \fill[white] (36.60,25.50) rectangle (84.60,68.60);
 \node[anchor=east,inner sep=0pt,font=\fontsize{6.0}{6.4}\selectfont] at (84.30,28.80) {核回归};
 \node[anchor=east,inner sep=0pt,font=\fontsize{6.0}{6.4}\selectfont] at (84.30,34.85) {径向 SVM};
 \node[anchor=east,inner sep=0pt,font=\fontsize{6.0}{6.4}\selectfont] at (84.30,40.90) {GBM};
 \node[anchor=east,inner sep=0pt,font=\fontsize{6.0}{6.4}\selectfont\bfseries] at (84.30,46.95) {HAL};
 \node[anchor=east,inner sep=0pt,font=\fontsize{6.0}{6.4}\selectfont] at (84.30,53.00) {多项式 MARS};
 \node[anchor=east,inner sep=0pt,font=\fontsize{6.0}{6.4}\selectfont] at (84.30,59.05) {回归树};
 \node[anchor=east,inner sep=0pt,font=\fontsize{6.0}{6.4}\selectfont] at (84.30,65.10) {随机森林};
 \fill[white] (148.80,25.50) rectangle (196.80,68.60);
 \node[anchor=east,inner sep=0pt,font=\fontsize{6.0}{6.4}\selectfont] at (196.50,28.80) {径向 SVM};
 \node[anchor=east,inner sep=0pt,font=\fontsize{6.0}{6.4}\selectfont] at (196.50,34.85) {核回归};
 \node[anchor=east,inner sep=0pt,font=\fontsize{6.0}{6.4}\selectfont] at (196.50,40.90) {GBM};
 \node[anchor=east,inner sep=0pt,font=\fontsize{6.0}{6.4}\selectfont\bfseries] at (196.50,46.95) {HAL};
 \node[anchor=east,inner sep=0pt,font=\fontsize{6.0}{6.4}\selectfont] at (196.50,53.00) {多项式 MARS};
 \node[anchor=east,inner sep=0pt,font=\fontsize{6.0}{6.4}\selectfont] at (196.50,59.05) {随机森林};
 \node[anchor=east,inner sep=0pt,font=\fontsize{6.0}{6.4}\selectfont] at (196.50,65.10) {回归树};
 \fill[white] (261.00,25.50) rectangle (309.00,68.60);
 \node[anchor=east,inner sep=0pt,font=\fontsize{6.0}{6.4}\selectfont] at (308.70,28.80) {GBM};
 \node[anchor=east,inner sep=0pt,font=\fontsize{6.0}{6.4}\selectfont\bfseries] at (308.70,34.85) {HAL};
 \node[anchor=east,inner sep=0pt,font=\fontsize{6.0}{6.4}\selectfont] at (308.70,40.90) {径向 SVM};
 \node[anchor=east,inner sep=0pt,font=\fontsize{6.0}{6.4}\selectfont] at (308.70,46.95) {核回归};
 \node[anchor=east,inner sep=0pt,font=\fontsize{6.0}{6.4}\selectfont] at (308.70,53.00) {随机森林};
 \node[anchor=east,inner sep=0pt,font=\fontsize{6.0}{6.4}\selectfont] at (308.70,59.05) {多项式 MARS};
 \node[anchor=east,inner sep=0pt,font=\fontsize{6.0}{6.4}\selectfont] at (308.70,65.10) {回归树};
 \fill[white] (373.20,25.50) rectangle (421.20,68.60);
 \node[anchor=east,inner sep=0pt,font=\fontsize{6.0}{6.4}\selectfont] at (420.90,28.80) {径向 SVM};
 \node[anchor=east,inner sep=0pt,font=\fontsize{6.0}{6.4}\selectfont] at (420.90,34.85) {GBM};
 \node[anchor=east,inner sep=0pt,font=\fontsize{6.0}{6.4}\selectfont\bfseries] at (420.90,40.90) {HAL};
 \node[anchor=east,inner sep=0pt,font=\fontsize{6.0}{6.4}\selectfont] at (420.90,46.95) {核回归};
 \node[anchor=east,inner sep=0pt,font=\fontsize{6.0}{6.4}\selectfont] at (420.90,53.00) {多项式 MARS};
 \node[anchor=east,inner sep=0pt,font=\fontsize{6.0}{6.4}\selectfont] at (420.90,59.05) {随机森林};
 \node[anchor=east,inner sep=0pt,font=\fontsize{6.0}{6.4}\selectfont] at (420.90,65.10) {回归树};
 \fill[white] (36.60,91.40) rectangle (84.60,134.50);
 \node[anchor=east,inner sep=0pt,font=\fontsize{6.0}{6.4}\selectfont] at (84.30,94.70) {回归树};
 \node[anchor=east,inner sep=0pt,font=\fontsize{6.0}{6.4}\selectfont] at (84.30,100.75) {GBM};
 \node[anchor=east,inner sep=0pt,font=\fontsize{6.0}{6.4}\selectfont\bfseries] at (84.30,106.80) {HAL};
 \node[anchor=east,inner sep=0pt,font=\fontsize{6.0}{6.4}\selectfont] at (84.30,112.85) {核回归};
 \node[anchor=east,inner sep=0pt,font=\fontsize{6.0}{6.4}\selectfont] at (84.30,118.90) {径向 SVM};
 \node[anchor=east,inner sep=0pt,font=\fontsize{6.0}{6.4}\selectfont] at (84.30,124.95) {随机森林};
 \node[anchor=east,inner sep=0pt,font=\fontsize{6.0}{6.4}\selectfont] at (84.30,131.00) {多项式 MARS};
 \fill[white] (148.80,91.40) rectangle (196.80,134.50);
 \node[anchor=east,inner sep=0pt,font=\fontsize{6.0}{6.4}\selectfont\bfseries] at (196.50,94.70) {HAL};
 \node[anchor=east,inner sep=0pt,font=\fontsize{6.0}{6.4}\selectfont] at (196.50,100.75) {GBM};
 \node[anchor=east,inner sep=0pt,font=\fontsize{6.0}{6.4}\selectfont] at (196.50,106.80) {随机森林};
 \node[anchor=east,inner sep=0pt,font=\fontsize{6.0}{6.4}\selectfont] at (196.50,112.85) {回归树};
 \node[anchor=east,inner sep=0pt,font=\fontsize{6.0}{6.4}\selectfont] at (196.50,118.90) {核回归};
 \node[anchor=east,inner sep=0pt,font=\fontsize{6.0}{6.4}\selectfont] at (196.50,124.95) {径向 SVM};
 \node[anchor=east,inner sep=0pt,font=\fontsize{6.0}{6.4}\selectfont] at (196.50,131.00) {多项式 MARS};
 \fill[white] (261.00,91.40) rectangle (309.00,134.50);
 \node[anchor=east,inner sep=0pt,font=\fontsize{6.0}{6.4}\selectfont\bfseries] at (308.70,94.70) {HAL};
 \node[anchor=east,inner sep=0pt,font=\fontsize{6.0}{6.4}\selectfont] at (308.70,100.75) {GBM};
 \node[anchor=east,inner sep=0pt,font=\fontsize{6.0}{6.4}\selectfont] at (308.70,106.80) {随机森林};
 \node[anchor=east,inner sep=0pt,font=\fontsize{6.0}{6.4}\selectfont] at (308.70,112.85) {回归树};
 \node[anchor=east,inner sep=0pt,font=\fontsize{6.0}{6.4}\selectfont] at (308.70,118.90) {核回归};
 \node[anchor=east,inner sep=0pt,font=\fontsize{6.0}{6.4}\selectfont] at (308.70,124.95) {径向 SVM};
 \node[anchor=east,inner sep=0pt,font=\fontsize{6.0}{6.4}\selectfont] at (308.70,131.00) {多项式 MARS};
 \fill[white] (373.20,91.40) rectangle (421.20,134.50);
 \node[anchor=east,inner sep=0pt,font=\fontsize{6.0}{6.4}\selectfont\bfseries] at (420.90,94.70) {HAL};
 \node[anchor=east,inner sep=0pt,font=\fontsize{6.0}{6.4}\selectfont] at (420.90,100.75) {GBM};
 \node[anchor=east,inner sep=0pt,font=\fontsize{6.0}{6.4}\selectfont] at (420.90,106.80) {随机森林};
 \node[anchor=east,inner sep=0pt,font=\fontsize{6.0}{6.4}\selectfont] at (420.90,112.85) {回归树};
 \node[anchor=east,inner sep=0pt,font=\fontsize{6.0}{6.4}\selectfont] at (420.90,118.90) {核回归};
 \node[anchor=east,inner sep=0pt,font=\fontsize{6.0}{6.4}\selectfont] at (420.90,124.95) {径向 SVM};
 \node[anchor=east,inner sep=0pt,font=\fontsize{6.0}{6.4}\selectfont] at (420.90,131.00) {多项式 MARS};
 \fill[white] (36.60,156.20) rectangle (84.60,199.30);
 \node[anchor=east,inner sep=0pt,font=\fontsize{6.0}{6.4}\selectfont] at (84.30,159.50) {核回归};
 \node[anchor=east,inner sep=0pt,font=\fontsize{6.0}{6.4}\selectfont] at (84.30,165.55) {径向 SVM};
 \node[anchor=east,inner sep=0pt,font=\fontsize{6.0}{6.4}\selectfont] at (84.30,171.60) {GBM};
 \node[anchor=east,inner sep=0pt,font=\fontsize{6.0}{6.4}\selectfont\bfseries] at (84.30,177.65) {HAL};
 \node[anchor=east,inner sep=0pt,font=\fontsize{6.0}{6.4}\selectfont] at (84.30,183.70) {回归树};
 \node[anchor=east,inner sep=0pt,font=\fontsize{6.0}{6.4}\selectfont] at (84.30,189.75) {多项式 MARS};
 \node[anchor=east,inner sep=0pt,font=\fontsize{6.0}{6.4}\selectfont] at (84.30,195.80) {随机森林};
 \fill[white] (148.80,156.20) rectangle (196.80,199.30);
 \node[anchor=east,inner sep=0pt,font=\fontsize{6.0}{6.4}\selectfont] at (196.50,159.50) {核回归};
 \node[anchor=east,inner sep=0pt,font=\fontsize{6.0}{6.4}\selectfont] at (196.50,165.55) {随机森林};
 \node[anchor=east,inner sep=0pt,font=\fontsize{6.0}{6.4}\selectfont\bfseries] at (196.50,171.60) {HAL};
 \node[anchor=east,inner sep=0pt,font=\fontsize{6.0}{6.4}\selectfont] at (196.50,177.65) {回归树};
 \node[anchor=east,inner sep=0pt,font=\fontsize{6.0}{6.4}\selectfont] at (196.50,183.70) {径向 SVM};
 \node[anchor=east,inner sep=0pt,font=\fontsize{6.0}{6.4}\selectfont] at (196.50,189.75) {GBM};
 \node[anchor=east,inner sep=0pt,font=\fontsize{6.0}{6.4}\selectfont] at (196.50,195.80) {多项式 MARS};
 \fill[white] (261.00,156.20) rectangle (309.00,199.30);
 \node[anchor=east,inner sep=0pt,font=\fontsize{6.0}{6.4}\selectfont] at (308.70,159.50) {核回归};
 \node[anchor=east,inner sep=0pt,font=\fontsize{6.0}{6.4}\selectfont] at (308.70,165.55) {随机森林};
 \node[anchor=east,inner sep=0pt,font=\fontsize{6.0}{6.4}\selectfont\bfseries] at (308.70,171.60) {HAL};
 \node[anchor=east,inner sep=0pt,font=\fontsize{6.0}{6.4}\selectfont] at (308.70,177.65) {回归树};
 \node[anchor=east,inner sep=0pt,font=\fontsize{6.0}{6.4}\selectfont] at (308.70,183.70) {GBM};
 \node[anchor=east,inner sep=0pt,font=\fontsize{6.0}{6.4}\selectfont] at (308.70,189.75) {径向 SVM};
 \node[anchor=east,inner sep=0pt,font=\fontsize{6.0}{6.4}\selectfont] at (308.70,195.80) {多项式 MARS};
 \fill[white] (373.20,156.20) rectangle (421.20,199.30);
 \node[anchor=east,inner sep=0pt,font=\fontsize{6.0}{6.4}\selectfont] at (420.90,159.50) {核回归};
 \node[anchor=east,inner sep=0pt,font=\fontsize{6.0}{6.4}\selectfont] at (420.90,165.55) {随机森林};
 \node[anchor=east,inner sep=0pt,font=\fontsize{6.0}{6.4}\selectfont\bfseries] at (420.90,171.60) {HAL};
 \node[anchor=east,inner sep=0pt,font=\fontsize{6.0}{6.4}\selectfont] at (420.90,177.65) {回归树};
 \node[anchor=east,inner sep=0pt,font=\fontsize{6.0}{6.4}\selectfont] at (420.90,183.70) {GBM};
 \node[anchor=east,inner sep=0pt,font=\fontsize{6.0}{6.4}\selectfont] at (420.90,189.75) {径向 SVM};
 \node[anchor=east,inner sep=0pt,font=\fontsize{6.0}{6.4}\selectfont] at (420.90,195.80) {多项式 MARS};
 \fill[white] (36.60,222.10) rectangle (84.60,265.20);
 \node[anchor=east,inner sep=0pt,font=\fontsize{6.0}{6.4}\selectfont] at (84.30,225.40) {GBM};
 \node[anchor=east,inner sep=0pt,font=\fontsize{6.0}{6.4}\selectfont] at (84.30,231.45) {核回归};
 \node[anchor=east,inner sep=0pt,font=\fontsize{6.0}{6.4}\selectfont\bfseries] at (84.30,237.50) {HAL};
 \node[anchor=east,inner sep=0pt,font=\fontsize{6.0}{6.4}\selectfont] at (84.30,243.55) {回归树};
 \node[anchor=east,inner sep=0pt,font=\fontsize{6.0}{6.4}\selectfont] at (84.30,249.60) {径向 SVM};
 \node[anchor=east,inner sep=0pt,font=\fontsize{6.0}{6.4}\selectfont] at (84.30,255.65) {多项式 MARS};
 \node[anchor=east,inner sep=0pt,font=\fontsize{6.0}{6.4}\selectfont] at (84.30,261.70) {随机森林};
 \fill[white] (148.80,222.10) rectangle (196.80,265.20);
 \node[anchor=east,inner sep=0pt,font=\fontsize{6.0}{6.4}\selectfont\bfseries] at (196.50,225.40) {HAL};
 \node[anchor=east,inner sep=0pt,font=\fontsize{6.0}{6.4}\selectfont] at (196.50,231.45) {随机森林};
 \node[anchor=east,inner sep=0pt,font=\fontsize{6.0}{6.4}\selectfont] at (196.50,237.50) {核回归};
 \node[anchor=east,inner sep=0pt,font=\fontsize{6.0}{6.4}\selectfont] at (196.50,243.55) {回归树};
 \node[anchor=east,inner sep=0pt,font=\fontsize{6.0}{6.4}\selectfont] at (196.50,249.60) {GBM};
 \node[anchor=east,inner sep=0pt,font=\fontsize{6.0}{6.4}\selectfont] at (196.50,255.65) {径向 SVM};
 \node[anchor=east,inner sep=0pt,font=\fontsize{6.0}{6.4}\selectfont] at (196.50,261.70) {多项式 MARS};
 \fill[white] (261.00,222.10) rectangle (309.00,265.20);
 \node[anchor=east,inner sep=0pt,font=\fontsize{6.0}{6.4}\selectfont\bfseries] at (308.70,225.40) {HAL};
 \node[anchor=east,inner sep=0pt,font=\fontsize{6.0}{6.4}\selectfont] at (308.70,231.45) {随机森林};
 \node[anchor=east,inner sep=0pt,font=\fontsize{6.0}{6.4}\selectfont] at (308.70,237.50) {核回归};
 \node[anchor=east,inner sep=0pt,font=\fontsize{6.0}{6.4}\selectfont] at (308.70,243.55) {GBM};
 \node[anchor=east,inner sep=0pt,font=\fontsize{6.0}{6.4}\selectfont] at (308.70,249.60) {回归树};
 \node[anchor=east,inner sep=0pt,font=\fontsize{6.0}{6.4}\selectfont] at (308.70,255.65) {径向 SVM};
 \node[anchor=east,inner sep=0pt,font=\fontsize{6.0}{6.4}\selectfont] at (308.70,261.70) {多项式 MARS};
 \fill[white] (373.20,222.10) rectangle (421.20,265.20);
 \node[anchor=east,inner sep=0pt,font=\fontsize{6.0}{6.4}\selectfont\bfseries] at (420.90,225.40) {HAL};
 \node[anchor=east,inner sep=0pt,font=\fontsize{6.0}{6.4}\selectfont] at (420.90,231.45) {核回归};
 \node[anchor=east,inner sep=0pt,font=\fontsize{6.0}{6.4}\selectfont] at (420.90,237.50) {随机森林};
 \node[anchor=east,inner sep=0pt,font=\fontsize{6.0}{6.4}\selectfont] at (420.90,243.55) {GBM};
 \node[anchor=east,inner sep=0pt,font=\fontsize{6.0}{6.4}\selectfont] at (420.90,249.60) {回归树};
 \node[anchor=east,inner sep=0pt,font=\fontsize{6.0}{6.4}\selectfont] at (420.90,255.65) {径向 SVM};
 \node[anchor=east,inner sep=0pt,font=\fontsize{6.0}{6.4}\selectfont] at (420.90,261.70) {多项式 MARS};
\end{tikzpicture}%
}

 \caption{模拟研究结果}
 \label{fig:ch06-prediction-simulation}
\end{figure}
\end{landscape}
\FloatBarrier

虽然 HAL 在光滑型和正弦型数据生成机制下没有排名第一，但其表现与那些通常被视为先进方法、且表现最好的机器学习算法相当。

\subsection{预测数据分析}
\label{subsec:ch06-prediction-analysis}

我们分别分析了五个公开可用的数据集，表 6.1 列出了这些数据集及其参考文献。样本量从 201 到 654 不等，$d$ 从 4 到 11 不等。除了模拟中评价的非参数方法外，我们还考虑用几种参数方法估计 $\bar Q_0$。其中包括只含主效应项的广义线性模型（generalized linear model，GLM）、基于 AIC 且纳入二阶交互项的逐步 GLM，以及通过十折交叉验证确定样条阶数的广义可加模型（generalized additive model，GAM）。我们还纳入了超级学习器和离散超级学习器，并将这九种算法分别作为候选算法。

\begin{table}[htbp]
 \centering
 \caption{用 HAL 估计量及其他比较方法分析的数据集}
 \label{tab:ch06-datasets}
 \begin{tabular}{lrr}
  \toprule
  名称 & $n$ & $d$\\
  \midrule
  cpu（Kibler et al.，1989） & 209 & 6\\
  laheart（Afifi and Azen，1979） & 201 & 11\\
  oecdpanel（Liu and Stengos，1999） & 616 & 6\\
  pima（Smith et al.，1988） & 392 & 7\\
  fev（Rosner，1999） & 654 & 4\\
  \bottomrule
 \end{tabular}
\end{table}

为了比较不同方法在不同结局、不同数据集上的表现，我们考察了每种方法的十折交叉验证均方误差，相对于主效应项 GLM 的均方误差的比值。大于 1 的值表示 GLM 的表现更好。各数据分析的结果见图 6.4。灰色圆点表示某个数据集中的相对 MSE，黑色十字表示全部五项研究的几何均值。超级学习器和离散超级学习器表现最好，其次是 HAL 估计量。HAL 估计量在 cpu 数据集上表现尤其出色，其交叉验证 MSE 几乎只有主效应项 GLM 的一半。

\begin{figure}[htbp]
 \centering
 % 原始图中各数据点、参照线、刻度和顺序全部保留；仅替换文字标签。
\resizebox{.88\textwidth}{!}{%
\begin{tikzpicture}[x=1bp,y=-1bp]
 \node[anchor=north west,inner sep=0pt] at (0,0) {\includegraphics[width=292bp,height=188bp]{figures/ch06_fig4_source.pdf}};
 \fill[white] (0,0) rectangle (88,159);
 \fill[white] (169,175) rectangle (230,188);
 \node[font=\fontsize{8.2}{9}\selectfont] at (195.9,182.7) {相对均方误差};
 \node[anchor=east,inner sep=0pt,font=\fontsize{8.0}{9}\selectfont] at (86.5,151.475) {径向 SVM};
 \node[anchor=east,inner sep=0pt,font=\fontsize{8.0}{9}\selectfont] at (86.5,136.290) {回归树};
 \node[anchor=east,inner sep=0pt,font=\fontsize{8.0}{9}\selectfont] at (86.5,122.014) {多项式 MARS};
 \node[anchor=east,inner sep=0pt,font=\fontsize{8.0}{9}\selectfont] at (86.5,107.735) {逐步 GLM};
 \node[anchor=east,inner sep=0pt,font=\fontsize{8.0}{9}\selectfont] at (86.5,94.169) {GBM};
 \node[anchor=east,inner sep=0pt,font=\fontsize{8.0}{9}\selectfont] at (86.5,78.984) {核回归};
 \node[anchor=east,inner sep=0pt,font=\fontsize{8.0}{9}\selectfont] at (86.5,65.468) {随机森林};
 \node[anchor=east,inner sep=0pt,font=\fontsize{8.0}{9}\selectfont] at (86.5,51.189) {GAM};
 \node[anchor=east,inner sep=0pt,font=\fontsize{8.0}{9}\selectfont\bfseries] at (86.5,36.669) {HAL};
 \node[anchor=east,inner sep=0pt,font=\fontsize{8.0}{9}\selectfont] at (86.5,21.774) {离散超级学习器};
 \node[anchor=east,inner sep=0pt,font=\fontsize{8.0}{9}\selectfont] at (86.5,7.447) {超级学习器};
\end{tikzpicture}%
}

 \caption{不同方法在五个真实数据集中的相对交叉验证均方误差。圆圈表示单个数据集的结果，十字表示五个数据集的几何均值。}
 \label{fig:ch06-relative-mse}
\end{figure}
\FloatBarrier
